\documentclass{article}
\usepackage[T1]{fontenc}
\usepackage{lmodern}
\usepackage{graphicx}
\usepackage{amsmath,amssymb}
\usepackage[a4paper,margin=2cm]{geometry}
\usepackage[absolute,overlay]{textpos}
\setlength{\TPHorizModule}{1mm}
\setlength{\TPVertModule}{1mm}
\usepackage[indonesian]{babel}
\usepackage{hyperref}
\setlength{\emergencystretch}{2em}
% unit: Halaman 10

\begin{document}

% === Halaman 10 (llm) ===

\section*{Prosedur Pembangkitan Kunci}

\begin{enumerate}
    \item Pilih sembarang bilangan prima $p$ ($p$ dapat di-share di antara anggota kelompok)
    \item Pilih dua buah bilangan acak, $g$ dan $x$, dengan syarat $g < p$, $g$ akar primitif dari $p$, dan $2 \le x \le p - 2$
    \item Hitung $y = g^x \pmod{p}$ \hfill (1)
\end{enumerate}

\noindent Hasil dari algoritma ini:
\begin{itemize}
    \item Kunci publik: tripel $(y, g, p)$
    \item Kunci privat: pasangan $(x, p)$
\end{itemize}

\end{document}
